
% ---
% Inserir folha de aprovação
% ---

% Isto é um exemplo de Folha de aprovação, elemento obrigatório da NBR
% 14724/2011 (seção 4.2.1.3). Você pode utilizar este modelo até a aprovação
% do trabalho. Após isso, substitua todo o conteúdo deste arquivo por uma
% imagem da página assinada pela banca com o comando abaixo:
%
% \includepdf{folhadeaprovacao_final.pdf}
%
\begin{folhadeaprovacao}

    \begin{center}
        {\fontseries{b}\selectfont\MakeTextUppercase{\normalsize\imprimirautor}}
    \end{center}
    
    \vfill
    
    \begin{center}
        {\fontseries{b}\selectfont\MakeTextUppercase{\imprimirtitulo}}
    \end{center}
    
    \vfill
    
    \abntex@ifnotempty{\imprimirpreambulo}{%
        \hspace{.45\textwidth}
        \begin{minipage}{.5\textwidth}
            \SingleSpacing
            \imprimirpreambulo\par
            \vspace*{4pt}
            {\imprimirorientadorRotulo~\imprimirorientador\par}
            \abntex@ifnotempty{\imprimircoorientador}{%
                {\imprimircoorientadorRotulo~\imprimircoorientador}%
            }%
        \end{minipage}%
    }
    
    \vfill

    \vspace*{1cm}

    % Aprovada em: \shortstack{20 \\ \underline{\hspace{1cm}}} / \shortstack{03 \\ \underline{\hspace{1cm}}} / \shortstack{2026 \\ \underline{\hspace{1.5cm}}}.    
    % Conceito obtido: \shortstack{10 \\ \underline{\hspace{1.5cm}}}

    \noindent % Garante que não haja recuo de parágrafo no início da linha
    Data de Aprovação: \shortstack{20 \\ \underline{\hspace{1cm}}} / \shortstack{03 \\ \underline{\hspace{1cm}}} / \shortstack{2026 \\ \underline{\hspace{1.5cm}}}.
    \hfill 
    Nota: \shortstack{10 \\ \underline{\hspace{1.5cm}}}.
    
    \vspace*{0.5cm}

    \begin{center}
        {\fontseries{b}\selectfont BANCA EXAMINADORA:}
    \end{center}
    
    \vspace*{1cm}
    
    \begin{center}
        \underline{\hspace{12cm}} \\
        \textbf{Prof. Dr. Almir Souza e Silva Neto} (Orientador) \\
        Instituto Federal de Educação, Ciência e Tecnologia do Maranhão (IFMA)
        
        \vspace*{1cm}
        
        \underline{\hspace{12cm}} \\
        \textbf{Prof. Me. Francisco dos Santos Viana} (Coorientador) \\
        Instituto Federal de Educação, Ciência e Tecnologia do Maranhão (IFMA)
        
        \vspace*{1cm}
        
        \underline{\hspace{12cm}} \\
        \textbf{Prof. Dr. Ronaldo Ribeiro Correa} (Examinador 1) \\
        Instituto Federal de Educação, Ciência e Tecnologia do Maranhão (IFMA)
        
        \vspace*{1cm}
        
        \underline{\hspace{12cm}} \\
        \textbf{Prof. Esp. Airton Augusto dos Anjos Costa} (Examinador 2)\\
        Instituto Federal de Educação, Ciência e Tecnologia do Maranhão (IFMA)
    \end{center}
    
    % \vspace*{\fill}
    
    % \begin{center}
    %     \imprimirlocal, 20 de março de 2026
    % \end{center}

    
\end{folhadeaprovacao}


%\textbf{	{Orientador: \vspace{-16pt} }
%	\assinatura{\textbf{Prof. \imprimirorientador , Dr.} \\ Univ. XXX} 
%	{Coorientador: \vspace{-16pt}}   
%	\assinatura{\textbf{Prof. \imprimircoorientador , Dr.} \\ Univ. XXX}
%	
%	{Membros: \vspace{-16pt} } 
%	
%	% --- Exemplo de assinaturas em sequência ---       
%	\setlength{\ABNTEXsignwidth}{8.5cm}
%	
%	\assinatura{\textbf{Prof. Professor, Dr.} \\ Univ. XXX}
%	\assinatura{\textbf{Prof. Professor, Dr.} \\ Univ. XXX}
%	\assinatura{\textbf{Prof. Professor, Dr.} \\ Univ. XXX}
%	
%	% --- Exemplo de assinaturas lado a lado ---
%	\setlength{\ABNTEXsignwidth}{7.5cm}
	%
	%    \noindent\hfill\assinatura*{\textbf{Prof. Professor, Dr.} \\ Univ. XXX}%
	%    \hfill%
	%    \assinatura*{\textbf{Prof. Professor, Dr.} \\ Univ. XXX}%
	%    \hfill
	%    
	%    \noindent\hfill\assinatura*{\textbf{Prof. Professor, Dr.} \\ Univ. XXX}%
	%    \hfill%
	%    \assinatura*{\textbf{Prof. Professor, Dr.} \\ Univ. XXX}%
	%    \hfill}